\section*{Глава 1. Анализ предметной области}
\addcontentsline{toc}{section}{Глава 1. Анализ предметной области}
\stepcounter{section}

\subsection{Обзор систем имитационного моделирования и управления ирригацией}

На современном рынке программного обеспечения для сельского хозяйства
выделяются три подхода к решению задачи мониторинга и управления
ирригацией: наукоёмкие профессиональные агрономические модели,
коммерческие IoT-платформы с реальными датчиками, а также
развлекательные и обучающие симуляторы фермерского хозяйства.
Для выявления незанятой ниши проведён сравнительный анализ
существующих решений по ключевым критериям: точность математической
модели, наличие геймификации, порог вхождения для пользователей
и ориентированность на обучение персонала.

\subsubsection{Профессиональные агрономические симуляторы}

\textbf{AquaCrop}~\cite{aquacrop_fao} (ФАО ООН, текущая версия~7.1) ---
стандартизированная имитационная модель, разработанная Продовольственной
и сельскохозяйственной организацией Объединённых Наций. Модель
ориентирована прежде всего на симуляцию урожайности травянистых
культур (herbaceous crops) в зависимости от водного стресса, оперируя
понятиями эвапотранспирации, водного баланса почвы и фенологических
фаз развития~\cite{steduto2009_aquacrop}. Для плодовых деревьев и
виноградников ФАО выпускает отдельное руководство (Irrigation and
Drainage Paper~66, <<Yield response to water of fruit trees and vines>>),
требующее ручной адаптации параметров модели. Несмотря на высокую
научную точность, AquaCrop имеет устаревший десктопный интерфейс
(Windows), высокий порог вхождения и полностью лишён элементов
интерактивности в реальном времени. Целевая аудитория --- учёные-агрономы
и аналитики международных организаций, а не линейный персонал ферм.

\textbf{DSSAT}~\cite{dssat_net} (Decision Support System for
Agrotechnology Transfer, версия~4.8.5, декабрь~2024) --- комплекс
имитационного моделирования, ведущий историю с 1989~года и включающий
более 45~моделей роста культур семейств CERES, CROPGRO и
других~\cite{jones2003_dssat}. Используется преимущественно в
научно-исследовательских целях для прогнозирования урожайности при
различных сценариях полива и климата. Взаимодействие с системой
происходит через оконное приложение для Windows и сложные
конфигурационные файлы. Система не предназначена для обучения
линейного персонала ферм (операторов полива) ввиду перегруженности
интерфейса и отсутствия веб-доступа.

\textbf{APSIM}~\cite{apsim_info} (Agricultural Production Systems
sIMulator) --- австралийская open-source платформа моделирования,
разрабатываемая с начала 1990-х годов и переведённая в открытый код
в~2014~году под управлением некоммерческой организации
APSIM~Initiative~\cite{holzworth2014_apsim}. Содержит более~80~модулей,
описывающих процессы в системах <<почва--растение--атмосфера>>, и
отличается высокой модульностью. Позволяет симулировать сложные
севообороты, влияние микроклимата и водный баланс для широкого спектра
культур. Как и DSSAT, ориентирована на агрономов-аналитиков и
академическую среду, а не на повседневный мониторинг и интерактивное
обучение персонала.

Общий недостаток всех трёх систем --- отсутствие веб-доступа,
невозможность совместной работы нескольких пользователей с различными
ролями и полное отсутствие вовлекающих механик, способных мотивировать
низкоквалифицированных операторов к регулярному взаимодействию
с системой.

\subsubsection{Коммерческие IoT-платформы как контекст для симуляции}

Параллельно с моделирующим ПО существуют платформы, работающие с
реальными датчиками. \textbf{Semios}~\cite{semios_official,semios_aginvesting}
(Канада) предлагает комплексную платформу <<Precision Agriculture
as a Service>> с mesh-сетью датчиков и подпиской 60--260~\$ за акр в
год. \textbf{CropX}~\cite{cropx_official,cropx_thespoon} (Израиль/США) ---
облачная аналитика влажности почвы с собственными винтовыми датчиками
и платной подпиской. \textbf{Arable}~\cite{arable_mark3_official,arable_futurefarming}
(США, спин-офф Принстонского университета) поставляет автономные
солнечные датчики Mark~3, измеряющие свыше сорока параметров (погода,
NDVI, канопи-температура). \textbf{Hortau}~\cite{hortau_official,hortau_growingproduce}
(Калифорния / Квебек) применяет подход тензиометрии почвы вместо
измерения объёмной влажности. Открытые решения представлены проектами
\textbf{LiteFarm}~\cite{litefarm_official,litefarm_github}
(React/Node.js, Университет Британской Колумбии) и
\textbf{Tania}~\cite{tania_github,tania_usetania_site} (Go/Vue.js,
Индонезия, с заявленной поддержкой MQTT).

Все перечисленные платформы либо ориентированы на крупные
агрохолдинги с высокой стоимостью входа, либо требуют приобретения
физического оборудования. Для малых фермерских хозяйств, не
располагающих бюджетом на IoT-инфраструктуру, \textbf{имитационная
модель} становится единственным доступным инструментом обучения
персонала и отработки агротехнических сценариев без риска потери
реального урожая и без капитальных затрат на сеть датчиков.

\subsubsection{<<Серьёзные игры>> и фермерские симуляторы}

\textbf{Farming Simulator}~\cite{farming_simulator_official} (GIANTS
Software, Швейцария) --- наиболее популярный развлекательный продукт
в агросфере: по состоянию на 2025~год суммарные продажи серии превысили
40~миллионов копий, а последняя часть (Farming Simulator~25) за первый
год реализована тиражом свыше четырёх
миллионов~\cite{giants_fs25_4million}. Обладает высоким уровнем
геймификации и фотореалистичной графикой, однако фокус сделан на
симуляцию управления техникой и экономику хозяйства. Биологические
процессы (влажность почвы, стресс растений, точечная ирригация)
смоделированы крайне поверхностно и не имеют научной ценности.
Исследование Pavlenko et~al.~(2024)~\cite{pavlenko_2024_sciencedirect,pavlenko_2024_centaur_reading}
на примере DLC <<Precision Farming>> подтвердило, что Farming Simulator
эффективен для повышения общей осведомлённости о точном земледелии,
но не подходит для отработки конкретных агротехнических навыков.

\textbf{Учебные платформы управления водными ресурсами.} В академической
среде существуют пошаговые симуляторы, позволяющие принимать решения
о поливе в условных сценариях. Европейский проект
\textbf{SmartH2O}~\cite{smarth2o_cordis} (7~рамочная программа ЕС,
грант~619172, 2014--2017, координатор --- SUPSI совместно с Politecnico
di Milano) применяет очки, бейджи и таблицы лидеров совместно с
реальными IoT-счётчиками воды; в 2018~году проект получил премию
Digital Water Prize конференции Water Innovation
Europe~\cite{smarth2o_digital_water_prize_ec}. Однако SmartH2O
ориентирован на \emph{бытовое} водопотребление в городах Швейцарии,
Великобритании и Испании, а не на сельскохозяйственную ирригацию.
Большинство подобных проектов реализованы в виде настольных игр или
устаревших веб-приложений без потокового обновления состояния в
реальном времени.

\textbf{Потребительские устройства для комнатных растений.} На рынке
появляются устройства с механикой виртуального питомца (Тамагочи) для
ухода за комнатными растениями. Характерный пример ---
\textbf{Senso}~\cite{senso_solidtech_prelaunch,senso_engadget_ces2026}
(компания SolidTech, представлен на CES~2026): пиксельный виртуальный
питомец с ежедневными миссиями по уходу, привязанный к одному
горшечному растению через магнитный сенсор. Подобные продукты успешно
используют геймификацию и положительно оцениваются пользователями,
однако работают исключительно с одиночными растениями и не
масштабируются на уровень сельскохозяйственных секторов и плантаций.

\subsubsection{Сводный сравнительный анализ}

Результаты анализа систематизированы в таблице~\ref{tab:competitors}.

\begin{longtable}{|p{2.6cm}|p{2.4cm}|c|c|c|p{2cm}|}
	\caption{Сравнительный анализ имитационных платформ и симуляторов}
	\label{tab:competitors} \\
	\hline
	\textbf{Платформа} & \textbf{Аудитория} & \textbf{Модель} &
	\textbf{Геймиф.} & \textbf{Веб} & \textbf{Порог входа} \\ \hline
	\endfirsthead
	\hline
	\textbf{Платформа} & \textbf{Аудитория} & \textbf{Модель} &
	\textbf{Геймиф.} & \textbf{Веб} & \textbf{Порог входа} \\ \hline
	\endhead
	\hline
	\endfoot

	AquaCrop (ФАО)    & Учёные           & Высокая & --   & --   & Очень высокий \\ \hline
	DSSAT             & Учёные           & Высокая & --   & --   & Очень высокий \\ \hline
	APSIM             & Учёные           & Высокая & --   & +/-- & Высокий       \\ \hline
	Farming Simulator & Геймеры          & Низкая  & +    & --   & Низкий        \\ \hline
	SmartH2O          & Быт. потребители & Средняя & +    & +    & Средний       \\ \hline
	Senso и аналоги   & Домовладельцы    & Нет     & +    & +    & Низкий        \\ \hline
	Semios / CropX    & Крупные хозяйства & Высокая & --  & +    & Высокий       \\ \hline
	LiteFarm / Tania  & Малые фермы      & Низкая  & --   & +    & Средний       \\ \hline
	\textbf{Проект}   & \textbf{Агрономы, операторы} & \textbf{Средняя} & \textbf{+} & \textbf{+} & \textbf{Низкий} \\ \hline
\end{longtable}

\indent Условные обозначения: Модель --- сложность математического расчёта
агрономических процессов; Геймиф. --- наличие игровых механик
(индексы здоровья, анимации, система уведомлений); Веб --- работа в
браузере без установки тяжёлого ПО; <<+/-->> --- частичная поддержка.

Анализ показал, что \textbf{существует выраженный разрыв} между
строгими научными симуляторами, лишёнными удобного интерфейса и
геймификации, и игровыми продуктами, лишёнными адекватной
математической модели. Коммерческие IoT-платформы решают задачу
мониторинга, но требуют значительных капиталовложений в оборудование.
\textbf{Искомая ниша --- пересечение двух множеств}: имитационная
модель с физически осмысленной динамикой (эвапотранспирация, водный
баланс, индекс здоровья) и геймифицированный веб-интерфейс (виртуальная
лейка, прогресс-бары, анимация обратной связи, уведомления). Разработка
такого веб-приложения позволит эффективно обучать операторов полива и
отрабатывать агротехнические сценарии без риска потери урожая и затрат
на физические датчики.

\subsection{Анализ программных инструментов разработки}

Для реализации клиент-серверного имитационного приложения требуется
стек технологий, способный обеспечить параллельный расчёт состояний
множества секторов, мгновенную доставку изменений на клиентские
устройства и компактное развёртывание одним разработчиком.
Выбор каждого компонента сопровождался сравнительным анализом
альтернатив.

\subsubsection{Серверная платформа и движок симуляции}

В качестве языка серверной части рассматривались \textbf{Node.js}
(JavaScript/TypeScript) и \textbf{Go~(Golang)}~\cite{go_dev_official,effective_go}.
Node.js предоставляет единый язык на клиенте и сервере, но его модель
конкурентности основана на однопоточном цикле событий, что усложняет
реализацию параллельного пересчёта десятков секторов в фоновом режиме.
Go, напротив, использует легковесные потоки (\emph{горутины}) и
встроенный планировщик, что позволяет тривиально запустить
\emph{Simulation Engine} в отдельной горутине одновременно с обработкой
HTTP-запросов и WebSocket-соединений. Дополнительное преимущество
Go --- статическая типизация и компиляция в единый бинарный файл,
упрощающая контейнеризацию. В качестве HTTP-маршрутизатора выбрана
библиотека \texttt{go-chi}~\cite{go_chi_github,go_chi_pkgdev} ---
идиоматичный минималистичный роутер с поддержкой группировки
маршрутов и цепочек middleware, используемых для JWT-аутентификации и
проверки ролей.

\subsubsection{Обмен данными в реальном времени}

Для трансляции <<тиков>> имитационной модели на интерфейс оператора
сравнивались три подхода: периодический опрос сервера
(\emph{polling}), \emph{server-sent events}~(SSE) и протокол
\textbf{WebSocket} (RFC~6455)~\cite{rfc6455_ietf,mdn_websocket}.
Polling создаёт избыточную нагрузку на сервер и задерживает
уведомления; SSE поддерживает только однонаправленную передачу от
сервера к клиенту. WebSocket обеспечивает полнодуплексный канал с
минимальными накладными расходами, что критично для сценариев, когда
клиент инициирует полив, а сервер немедленно рассылает обновлённое
состояние сектора всем подключённым операторам и агрономам. В качестве
реализации выбрана библиотека
\texttt{gorilla/websocket}~\cite{gorilla_websocket_github,gorilla_websocket_pkgdev}
--- фактический стандарт в экосистеме Go.

\subsubsection{Клиентский фреймворк}

Для реализации клиентской части сравнивались
\textbf{React.js}~\cite{react_official,react_reference_hooks} и
\textbf{Vue.js}. Vue предлагает более пологую кривую обучения и
компактный шаблонный синтаксис, однако React обладает более развитой
экосистемой библиотек для анимации пользовательского интерфейса и
визуализации состояний, что важно для реализации геймифицированных
элементов --- анимации виртуальной лейки, прогресс-баров суточных
лимитов воды, плавных переходов между состояниями растения.
В качестве сборщика выбран \textbf{Vite}~\cite{vite_official,vite_guide}
--- современная альтернатива Webpack с мгновенной горячей перезагрузкой
модулей (HMR), что существенно ускоряет итеративную разработку
интерактивных компонентов.

\subsubsection{База данных и объектно-реляционное отображение}

В качестве СУБД выбрана
\textbf{PostgreSQL}~16~\cite{postgresql_official,postgresql_16_release_notes}.
Реляционная модель с поддержкой типов \texttt{UUID}, \texttt{TIMESTAMPTZ}
и \texttt{FLOAT} обеспечивает эффективное хранение конфигураций
секторов, логов полива и временных рядов телеметрии; составной индекс
по полям \texttt{(sector\_id, recorded\_at~DESC)} позволяет выбирать
историю показаний сектора за произвольный интервал с предсказуемой
производительностью. Взаимодействие с БД осуществляется через
ORM-библиотеку \textbf{Bun}~\cite{bun_uptrace_official,bun_github} ---
лёгкую альтернативу GORM, предоставляющую типобезопасный fluent-builder
запросов без скрытой <<магии>> и существенного оверхеда.

\subsubsection{Безопасность и ролевая модель}

Аутентификация реализована по схеме \textbf{JSON~Web~Token}
(RFC~7519)~\cite{rfc7519_ietf,jwt_io_intro} с подписью HMAC-SHA256:
сервер выдаёт клиенту токен со сроком жизни 24~часа, содержащий
идентификатор пользователя и роль (\texttt{agronomist} или
\texttt{operator}). Для хеширования паролей применяется алгоритм
\textbf{bcrypt}~\cite{provos_mazieres_1999,pkggo_bcrypt} со стандартным
коэффициентом стоимости. Middleware-цепочка верифицирует токен при
каждом запросе и инжектирует пользовательский контекст в дальнейшую
обработку.

\subsubsection{Контейнеризация}

Развёртывание приложения реализовано средствами \textbf{Docker
Compose}~\cite{docker_compose_docs,docker_compose_file_reference},
объединяющего три сервиса: базу данных PostgreSQL, API-сервер на Go и
фронтенд-приложение на Vite. Это обеспечивает воспроизводимость
окружения на любой машине с установленным Docker и упрощает
демонстрацию проекта.

Использование связки Go и React позволяет создать архитектуру, в
которой математическая логика имитационной модели защищена на сервере
и выполняется автономно независимо от активности клиентов, а
клиентское приложение занимается исключительно плавной отрисовкой
результатов симуляции и обработкой пользовательских действий.

\subsection{Формулировка цели и задач работы}

На основе проведённого анализа предметной области (см.
подраздел~1.1), сравнения существующих симуляторов и обоснования
архитектурного стека (см. подраздел~1.2) была сформулирована цель
исследования.

\textbf{Цель работы} --- разработка интерактивного веб-приложения для
мониторинга и управления ирригацией плантаций тропических культур
(на примере личи, \emph{Litchi chinensis}) на основе серверной
имитационной модели с элементами геймификации, ориентированного
на операторов полива и агрономов малых фермерских хозяйств.

Для достижения поставленной цели необходимо решить следующие
\textbf{задачи}:
\begin{enumerate}
	\item Провести анализ существующих имитационных моделей, обучающих
	игровых систем и коммерческих IoT-платформ в области
	управления ирригацией; выявить незанятую нишу на
	пересечении имитационного моделирования и геймификации.
	\item Спроектировать структуру реляционной базы данных для
	хранения конфигураций секторов, назначений операторов,
	логов полива, временных рядов телеметрии и
	пользовательских уведомлений.
	\item Разработать фоновый движок симуляции на языке Go,
	математически описывающий процессы изменения влажности
	почвы, дрейфа температуры воздуха и динамики индекса
	здоровья растений в ответ на засуху, переувлажнение и
	своевременный полив.
	\item Разработать клиентское приложение на React с ролевой
	моделью доступа (агроном / оператор) и геймифицированным
	интерфейсом управления виртуальным поливом: анимированная
	лейка, индексы здоровья, суточные лимиты воды,
	прогресс-бары.
	\item Интегрировать полнодуплексную связь по протоколу
	WebSocket для потоковой передачи изменений состояния
	модели на экран пользователя в реальном времени, включая
	немедленную доставку уведомлений о критических событиях
	(засуха, переувлажнение, падение индекса здоровья).
	\item Реализовать импорт конфигурации секторов из CSV и экспорт
	накопленной телеметрии для офлайн-анализа.
	\item Провести тестирование разработанной системы в условиях
	имитации реальной нагрузки на сервер и контейнерного
	развёртывания.
\end{enumerate}